\documentclass[12pt]{article}
\usepackage[T1]{fontenc}
\usepackage[utf8]{inputenc}
\usepackage{lmodern}
\usepackage{textcomp}
\usepackage{enumitem}
\usepackage{geometry}
\geometry{margin=1in}
\begin{document}
\begin{center}
Answer Key - History - Graduate Level Assessment 23 \
\small (Answers are ordered exactly as in the question paper.)
\end{center}

\section*{Multiple Choice}
\begin{enumerate}[label=\arabic*.]
\item \textbf{Answer:} B
\\ \textit{Options:} A) coins recovered from ancient shipwrecks.; B) dendrochronology.; C) all of the above.; D) ice core sampling.; E) written records.
\\ \textit{Note:} Dendrochronology provides precise year-by-year data through tree-ring analysis, allowing for the calibration of radiocarbon dating results against known chronological events, thus enhancing the accuracy of the dating process.
\item \textbf{Answer:} D
\\ \textit{Options:} A) pottery fragments.; B) irrigation canals.; C) jewelry artifacts.; D) monumental works.; E) pyramids.
\\ \textit{Note:} Monumental works serve as tangible representations of the power, resources, and cultural values of ancient state societies, reflecting their organizational complexity and societal achievements.
\item \textbf{Answer:} A
\\ \textit{Options:} A) a series of decisive military conquests against a rival king.; B) a depiction of him as a corn god, a symbol of fertility and prosperity.; C) a scene of him hunting and defeating a large jaguar.; D) an eagle, a serpent, and a jaguar, all of which were important royal animals.; E) him receiving a crown from the Mayan god of death and rebirth.
\\ \textit{Note:} The depiction of a series of decisive military conquests against a rival king on the stele symbolizes Yax Pahsaj's authority and legitimacy as a ruler in Mayan society, reflecting the importance of military success in establishing and maintaining royal power.
\item \textbf{Answer:} A
\\ \textit{Options:} A) the foramen magnum; B) the skull; C) the feet; D) the femur; E) the spine
\\ \textit{Note:} The foramen magnum's position at the base of the skull indicates whether the spine connects vertically (bipedalism) or horizontally (quadrupedalism) to the skull, thus determining the locomotion type in primates.
\item \textbf{Answer:} A
\\ \textit{Options:} A) 1,000 B.P.; B) 8,000 B.P.; C) 7,000 B.P.; D) 3,000 B.P.; E) 6,000 B.P.
\\ \textit{Note:} The Yang-shao culture transitioned to the Lung-Shan culture around 1,000 B.P. as a result of evolving agricultural practices and societal changes during the Neolithic period in China.
\end{enumerate}

\section*{True / False}
\begin{enumerate}[label=\arabic*.]
\setcounter{enumi}{5}
\item \textbf{Answer:} False
\\ \textit{Options:} A) True; B) False
\\ \textit{Note:} The correct answer is: oracles-principles or mechanisms
\item \textbf{Answer:} False
\\ \textit{Options:} A) True; B) False
\\ \textit{Note:} The correct answer is: 1860
\item \textbf{Answer:} True
\\ \textit{Options:} A) True; B) False
\\ \textit{Note:} According to the passage: 1883
\item \textbf{Answer:} True
\\ \textit{Options:} A) True; B) False
\\ \textit{Note:} According to the passage: Allied aircraft
\item \textbf{Answer:} False
\\ \textit{Options:} A) True; B) False
\\ \textit{Note:} The correct answer is: budget deficit
\end{enumerate}

\section*{Long Form Response}
\begin{enumerate}[label=\arabic*.]
\setcounter{enumi}{10}
\item \textbf{Answer:} Flowers
\\ \textit{Note:} Expected answer: Flowers
\item \textbf{Answer:} similarities to wool
\\ \textit{Note:} Expected answer: similarities to wool
\end{enumerate}

\vfill
\noindent\textit{End of Answer Key}
\end{document}